\documentclass[11pt]{article}
\usepackage{fontspec,amsmath,amssymb,amsthm,hyperref}
\usepackage[margin=1in]{geometry}
\setlength{\emergencystretch}{3em}
\title{The OBABO splitting and its exact pathwise identities}
\author{LangevinLean mathematical exposition}
\date{}
\begin{document}\maketitle
\noindent Source module: \texttt{KineticLangevin/OBABO.lean}. Compile this proof document with XeLaTeX or LuaLaTeX. The accompanying HTML pairs the mathematical argument with the exact Lean statements and records the scope of the formalization.


OBABO brackets a kick–drift–kick Hamiltonian step between two Ornstein–Uhlenbeck half steps. The module proves the composition and its reversal and cancellation identities for arbitrary supplied vectors. The Gaussian-law and convergence analysis is separate. Its thermostat coefficient is chosen so that the damping and fresh noise exactly preserve a target scalar velocity variance.


\section{Define each subflow and compose O–B–A–B–O}


In velocity–position ordering, \(O(v,x)=(\delta v+\sigma\xi,x)\), \(B_h(v,x)=(v-hg(x),x)\), and \(A_h(v,x)=(v,x+hv)\). \texttt{obaboStep} composes \(O,B_{h/2},A_h,B_{h/2},O\), using distinct noise arguments in the two thermostat steps. Expanding the nested definitions gives \(v_{1/2}=\delta v+\sigma\xi-hg(x)/2\), \(x^+=x+hv_{1/2}\), and \(v^+=\delta(v_{1/2}-hg(x^+)/2)+\sigma\xi'\). \texttt{obaboStep\_expanded} is proved by \texttt{rfl} because this formula follows directly from the composition; the new-position force is structurally enforced.



\noindent Lean declarations: \texttt{ornsteinStep}, \texttt{kickStep}, \texttt{driftStep}, \texttt{obaboStep}, \texttt{obaboStep\_expanded}.

\section{Reverse the deterministic kick and drift}


The kick leaves position fixed, so subtracting \(hg(x)\) and then adding it cancels exactly. The drift leaves velocity fixed, so adding \(hv\) and then subtracting it also cancels. Their inverse formulas hold for arbitrary real \(h\) and arbitrary force, with no smoothness assumption. \texttt{reverseVelocity} negates velocity and keeps position; applying it twice is the identity. Conjugating a kick or drift by this reversal changes \(h\) to \(-h\). The proofs use pair extensionality and simplification of scalar signs, making the time-reversal algebra explicit.



\noindent Lean declarations: \texttt{kickStep\_neg}, \texttt{driftStep\_neg}, \texttt{reverseVelocity}, \texttt{reverseVelocity\_twice}, \texttt{reverseVelocity\_kick}, \texttt{reverseVelocity\_drift}.

\section{The palindromic Hamiltonian composition has an exact inverse}


\texttt{verletStep} is \(B_{h/2}A_hB_{h/2}\). To invert it, reverse the order and invert each factor. Since the composition is palindromic, this is precisely \(B_{-h/2}A_{-h}B_{-h/2}\). Lean expands the two compositions, rewrites \((-h)/2=-(h/2)\), and cancels the three adjacent inverse pairs with the preceding lemmas. The OBABO zero-time theorem separately checks that \(h=0,\delta=1,\sigma=0\) makes every substep the identity. Reversal of the deterministic map alone does not assert reversibility of a stochastic transition with a particular probability density.



\noindent Lean declarations: \texttt{verletStep}, \texttt{verletStep\_neg}, \texttt{obaboStep\_zero}.

\section{Synchronous noises cancel while the new forces remain}


Apply two OBABO steps with the same \(\xi,\xi'\). The difference of half velocities is \(a=\delta\Delta v-h(g(x)-g(y))/2\). Thus \(\Delta x^+=\Delta x+ha\). The final shared refresh cancels too, leaving \(\Delta v^+=\delta(a-h(g(x^+)-g(y^+))/2)\). Lean unfolds all subflows, proves the velocity and position identities as a conjunction, and normalizes vector expressions with \texttt{module}. Although additive noises disappear explicitly from the difference, the new positions still depend on the common first noise. The identity preserves their force evaluations instead of treating them as old-position gradients.



\noindent Lean declarations: \texttt{obaboStep\_synchronous}.

\section{Choose damping and prove the thermal variance identity}


The OU half-step damping is \(\delta=e^{-\gamma h/2}>0\). The exponential addition identity gives \(\delta^2=e^{-\gamma h}\). With \(\gamma,h\ge0\), this is at most one, so \(s^2=1-\delta^2\ge0\). For a target velocity variance \(u\ge0\), set \(\sigma=\sqrt{u(1-\delta^2)}\). Then \(\delta^2u+\sigma^2=u\). Lean proves positivity directly from \texttt{exp\_pos}, rewrites the square with \texttt{exp\_add}, uses the nonpositive exponent to bound the exponential, and finally applies \texttt{Real.sq\_sqrt} followed by \texttt{ring}. For independent centered input velocity and fresh Gaussian noise this is the scalar covariance balance; Gaussian-law preservation additionally uses those distributional facts.



\noindent Lean declarations: \texttt{obaboDamping}, \texttt{obaboVariance}, \texttt{obaboDamping\_pos}, \texttt{obaboDamping\_sq}, \texttt{obaboVariance\_nonneg}, \texttt{obabo\_thermostat\_variance}.

\end{document}
